\section{Research questions and hypotheses}
\label{sec:rq}

We pre-registered the hypotheses before inspecting the full results (repository
file \texttt{research/EXPERIMENT\_SPEC.md}).

\begin{itemize}
  \item \textbf{RQ1.} Under matched conditions, does Azure-domain training reduce
  token count on unseen Azure services relative to general-English training?
  \item \textbf{RQ2.} What general-English regression accompanies full
  specialization?
  \item \textbf{RQ3.} Does a 50/50 Mixed corpus preserve most of the in-domain
  gain while limiting the regression?
  \item \textbf{RQ4.} Does the effect transfer to other technical-documentation
  domains?
  \item \textbf{RQ5.} Is the effect stable across 8K, 16K, and 32K vocabularies?
  \item \textbf{RQ6.} Does the effect persist on prose, rather than being explained
  only by code and markup?
  \item \textbf{RQ7.} How do compression and CPU tokenizer throughput relate?
\end{itemize}

\paragraph{Primary hypothesis.} At a fixed byte-level BPE algorithm, vocabulary
size, training-byte budget, normalization, sampling, and evaluation set, an
Azure-trained tokenizer produces fewer tokens per 1{,}000 UTF-8 bytes on the Azure
cross-service holdout than a byte-matched General tokenizer.

\paragraph{Null hypothesis.} After those controls, the Azure- and General-trained
tokenizers do not differ in token count on the held-out set, i.e.\ the point
estimate of the relative reduction is $\le 0$ or its 95\% paired-bootstrap CI is
compatible with no positive reduction.
